\documentclass{apcalc}
\title{Unit 9 Practice Set: Parametric Equations, Polar Coordinates, and Vector-Valued Functions}
\author{Paul Wang}
\date{November 2026}
\volume{Practice Set 9 of 10 (topics 9.1–9.9)}
\shorttitle{Unit 9 --- Practice Set}
\begin{document}
\maketitle
\begin{abstract}
Fifteen exam-style multiple-choice questions on Unit 9, grouped by the sections of the Unit 9 Study Guide, and one nine-point calculator-active free-response question on a particle moving in the plane. Items marked [calculator] allow a graphing calculator; the rest are no-calculator. Write each answer in the blank beside the number; the key on the last pages gives an answers-only strip, a reason and rule for every item, and the free-response scoring guide.
\end{abstract}
\scoringline{15 MC \quad\blank\ / 9 FR}{Exam pace: about 2 minutes per multiple-choice item (30 min) and 15 minutes for the free response.}
\section*{1\quad Parametric curves and their derivatives}
\begin{setrule}R9.1–R9.2. $\frac{dy}{dx}=\frac{y'(t)}{x'(t)}$; $\frac{d^2y}{dx^2}=\frac{d}{dt}\!\left(\frac{dy}{dx}\right)\big/x'(t)$.\end{setrule}
\begin{q}{1}
For the curve $x=t^2+1$, $y=t^3-t$, the value of $\dfrac{dy}{dx}$ at $t=2$ is
\chfour{$\dfrac{4}{11}$}{$\dfrac{11}{4}$}{$6$}{$11$}
\end{q}
\begin{q}{2}
At what point does the curve $x=t^3-3t$, $y=t^2-4t$ have a horizontal tangent line?
\chfour{$(2,-4)$}{$(-2,-3)$}{$(2,5)$}{$(0,0)$}
\end{q}
\begin{q}{3}
For $x=t^2+t$, $y=t^3$, the value of $\dfrac{d^2y}{dx^2}$ at $t=1$ is
\chfour{$\dfrac49$}{$1$}{$\dfrac43$}{$3$}
\end{q}
\section*{2\quad Parametric arc length}
\begin{setrule}R9.3. $L=\int\sqrt{x'(t)^2+y'(t)^2}\,dt$.\end{setrule}
\begin{q}{4}
What is the length of the path $x=t^2$, $y=\frac23t^3$ for $0\le t\le1$?
\chfour{$\dfrac{\sqrt{13}}{3}$}{$\dfrac23\left(2\sqrt2-1\right)$}{$\dfrac53$}{$\dfrac{4\sqrt2}{3}$}
\end{q}
\section*{3\quad Vector-valued functions and planar motion}
\begin{setrule}R9.4–R9.5. Speed $\sqrt{x'^2+y'^2}$; position = initial position + $\int$ velocity, componentwise.\end{setrule}
\begin{q}{5}
A particle has velocity $\langle3t^2,\,4t\rangle$. What is its speed at $t=1$?
\chfour{$\sqrt7$}{$5$}{$7$}{$25$}
\end{q}
\begin{q}{6}
A particle's position is $\langle\sin2t,\,t^3\rangle$. Its acceleration vector at $t=\frac\pi4$ is
\chtwo{$\left\langle0,\,\frac{3\pi^2}{16}\right\rangle$}{$\left\langle4,\,\frac{3\pi}{2}\right\rangle$}{$\left\langle-4,\,\frac{3\pi^2}{16}\right\rangle$}{$\left\langle-4,\,\frac{3\pi}{2}\right\rangle$}
\end{q}
\begin{q}{7}
A particle has velocity $\langle2t,\,e^t\rangle$ and is at $(1,3)$ when $t=0$. Where is it at $t=1$?
\chfour{$(1,e)$}{$(3,e+2)$}{$(2,e+2)$}{$(2,e+3)$}
\end{q}
\begin{q}{8}
A particle moves with velocity $\langle\cos(t^2),\,\sin t\rangle$. What is the total distance it travels for $0\le t\le2$?\calcmark
\chfour{$1.416$}{$1.489$}{$2.207$}{$2.472$}
\end{q}
\begin{q}{9}
A particle has velocity $\langle t^2-4,\,3-t\rangle$ for $0<t<5$. For which $t$ is it moving left and up?
\chfour{$0<t<2$}{$2<t<3$}{$3<t<5$}{$0<t<3$}
\end{q}
\section*{4\quad Polar curves and their derivatives}
\begin{setrule}R9.6–R9.8. $x=r\cos\theta$, $y=r\sin\theta$; distance from the origin grows when $r$ and $\frac{dr}{d\theta}$ share a sign.\end{setrule}
\begin{q}{10}
Which rectangular equation is equivalent to $r=6\sin\theta$?
\chfour{$x^2+y^2=6x$}{$y=6$}{$x^2+y^2=36$}{$x^2+y^2=6y$}
\end{q}
\begin{q}{11}
For the polar curve $r=1+\cos\theta$, the slope $\dfrac{dy}{dx}$ at $\theta=\frac\pi2$ is
\chfour{$-1$}{$0$}{$1$}{undefined}
\end{q}
\begin{q}{12}
For $r=2\cos(2\theta)$ at $\theta=\frac\pi3$, which statement is true?
\chone{The point moves toward the origin, because $\frac{dr}{d\theta}<0$.}{The point moves away from the origin, because $r$ and $\frac{dr}{d\theta}$ are both negative.}{The point moves toward the origin, because $r<0$.}{The point moves away from the origin, because $\frac{dr}{d\theta}>0$.}
\end{q}
\section*{5\quad Area in polar coordinates}
\begin{setrule}R9.9–R9.10. $A=\frac12\int r^2\,d\theta$; between curves $\frac12\int(r_{\text{out}}^2-r_{\text{in}}^2)\,d\theta$.\end{setrule}
\begin{q}{13}
The area of one petal of the rose $r=\sin3\theta$ is
\chfour{$\dfrac{\pi}{12}$}{$\dfrac{\pi}{6}$}{$\dfrac{\pi}{4}$}{$\dfrac{\pi}{3}$}
\end{q}
\begin{q}{14}
Which expression gives the area inside $r=3\cos\theta$ and outside $r=1+\cos\theta$?
\chone{$\displaystyle\frac12\int_{-\pi/3}^{\pi/3}(2\cos\theta-1)^2\,d\theta$}{$\displaystyle\frac12\int_{0}^{\pi/3}\left(9\cos^2\theta-(1+\cos\theta)^2\right)d\theta$}{$\displaystyle\frac12\int_{-\pi/3}^{\pi/3}\left(9\cos^2\theta-(1+\cos\theta)^2\right)d\theta$}{$\displaystyle\int_{-\pi/3}^{\pi/3}\left(9\cos^2\theta-(1+\cos\theta)^2\right)d\theta$}
\end{q}
\begin{q}{15}
What is the area of the region bounded by the polar curve $r=\theta+\sin\theta$ for $0\le\theta\le\pi$ (and the $x$-axis)?\calcmark
\chfour{$3.467$}{$5.168$}{$5.953$}{$9.095$}
\end{q}
\clearpage\section*{Free response}
Calculator required. Show the integral or expression you evaluate, and give decimal answers correct to three places.
\begin{frq}{Question 1 --- A particle moving in the plane (9 points, calculator required)}
A particle moves in the $xy$-plane so that $\dfrac{dx}{dt}=\ln(t^2+1)-1$ and $\dfrac{dy}{dt}=2\sin t$ for $t\ge0$. At $t=0$ the particle is at $(1,4)$.
\partlabel{a}Find the speed of the particle at $t=2$, and find its acceleration vector at $t=2$.\workspace{1.6in}
\partlabel{b}Find the $y$-coordinate of the particle's position at $t=3$.\workspace{1.6in}
\partlabel{c}Find the total distance traveled by the particle over $0\le t\le3$.\workspace{1.6in}
\partlabel{d}Find the position of the particle at $t=1$, and write an equation for the line tangent to its path at that point.\workspace{1.6in}
\end{frq}
\clearpage\section*{Answer key}
\keystrip{\keyitem{1}{B}\keyitem{2}{A}\keyitem{3}{A}\keyitem{4}{B}\keyitem{5}{B}\keyitem{6}{D}\keyitem{7}{C}\keyitem{8}{C}\keyitem{9}{A}\keyitem{10}{D}\keyitem{11}{C}\keyitem{12}{B}\keyitem{13}{A}\keyitem{14}{C}\keyitem{15}{D}}
\vspace{4pt}
\begin{keytable}
\keysection{1. Parametric curves and their derivatives}
1 & B & $\frac{3t^2-1}{2t}=\frac{11}{4}$. & R9.1\\
2 & A & $\frac{dy}{dt}=2t-4=0$ at $t=2$, where $\frac{dx}{dt}=9\neq0$. & R9.1\\
3 & A & $\frac{d}{dt}\frac{3t^2}{2t+1}=\frac{6t^2+6t}{(2t+1)^2}=\frac{12}{9}$; divide by $x'(1)=3$. & R9.2\\
\keysection{2. Parametric arc length}
4 & B & Speed $=\sqrt{4t^2+4t^4}=2t\sqrt{1+t^2}$; integrate. & R9.3\\
\keysection{3. Vector-valued functions and planar motion}
5 & B & $\sqrt{3^2+4^2}=5$. & R9.4\\
6 & D & $\langle-4\sin2t,\,6t\rangle$ at $\frac\pi4$. & R9.4\\
7 & C & $x=1+\int_0^12t\,dt=2$; $y=3+\int_0^1e^t\,dt=3+e-1$. & R9.5\\
8 & C & $\int_0^2\sqrt{\cos^2(t^2)+\sin^2t}\,dt\approx2.207$. & R9.5, \emph{calc}\\
9 & A & Left: $x'<0$ for $t<2$; up: $y'>0$ for $t<3$. & R9.4\\
\keysection{4. Polar curves and their derivatives}
10 & D & Multiply by $r$: $r^2=6r\sin\theta$. & R9.6\\
11 & C & $\frac{dy}{d\theta}=-1$ and $\frac{dx}{d\theta}=-1$ at $\frac\pi2$. & R9.7\\
12 & B & $r=-1$, $r'=-2\sqrt3$: same sign, so $|r|$ increases. & R9.8\\
\keysection{5. Area in polar coordinates}
13 & A & $\frac12\int_0^{\pi/3}\sin^23\theta\,d\theta=\frac{\pi}{12}$. & R9.9\\
14 & C & Intersections at $\cos\theta=\frac12$; outer minus inner, each squared; the area is $\pi$. & R9.10\\
15 & D & $\frac12\int_0^\pi(\theta+\sin\theta)^2\,d\theta\approx9.095$. & R9.9, \emph{calc}\\
\end{keytable}
\section*{Free-response scoring guide}
\subsection*{Question 1: A particle moving in the plane (9 points)}
\textbf{(a)} Speed $=\sqrt{x'(2)^2+y'(2)^2}=\sqrt{(\ln5-1)^2+(2\sin2)^2}\approx1.918$. $a(2)=\left\langle\frac{2t}{t^2+1},\,2\cos t\right\rangle_{t=2}=\langle0.8,\,-0.832\rangle$.
\begin{itemize}[leftmargin=1.6em,itemsep=0pt]
\item 1 point: speed $\approx1.918$
\item 1 point: acceleration $\langle0.8,\,-0.832\rangle$
\end{itemize}
\textbf{(b)} $y(3)=4+\int_0^32\sin t\,dt=4+2(1-\cos3)\approx7.980$.
\begin{itemize}[leftmargin=1.6em,itemsep=0pt]
\item 1 point: $4+\int_0^3 2\sin t\,dt$
\item 1 point: $y(3)\approx7.980$
\end{itemize}
\textbf{(c)} $\int_0^3\sqrt{(\ln(t^2+1)-1)^2+4\sin^2t}\,dt\approx4.786$.
\begin{itemize}[leftmargin=1.6em,itemsep=0pt]
\item 1 point: integral of speed with limits 0 and 3
\item 1 point: distance $\approx4.786$
\end{itemize}
\textbf{(d)} $x(1)=1+\int_0^1(\ln(t^2+1)-1)\,dt\approx0.264$; $y(1)=4+2(1-\cos1)\approx4.919$. $\frac{dy}{dx}=\frac{2\sin1}{\ln2-1}\approx-5.485$. Tangent: $y=4.919-5.485(x-0.264)$.
\begin{itemize}[leftmargin=1.6em,itemsep=0pt]
\item 1 point: $x(1)=1+\int_0^1x'(t)\,dt\approx0.264$ and $y(1)\approx4.919$
\item 1 point: slope $\frac{dy}{dx}\approx-5.485$
\item 1 point: tangent line $y=4.919-5.485(x-0.264)$
\end{itemize}
\end{document}